% GENERATED FILE. Edit canonical lessons, then run npm run generate:books.
% canonical-source-hash: fnv1a64:5972a1b96de2ca09
% canonical-lessons: TA-C331-aracanai, TA-C331-ceppetu, TA-C331-vallal, TA-C331-kotai, TA-C331-celi

\chapter{Decrees and Patrons}
\label{ch:ta-decrees-and-patrons}
\hlchaptermodality{331}

\begin{chapteropening}
\textbf{By the end of this chapter:} \emph{I can talk about royal orders and generous patrons.}

Complete the last lesson of chapter 331: I can talk about royal orders and generous patrons.
\end{chapteropening}

\section[\texorpdfstring{\ta{அரசாணை} (aracāṇai)}{அரசாணை (aracāṇai)}]{\ta{அரசாணை} (aracāṇai) — a royal decree}

\label{lesson:TA-C331-aracanai}



\begin{quote}
Before the new one: say the Tamil for a royal tribute, then the Tamil for to flourish.
\end{quote}

\subsection*{You'll want to know: \ta{அரசாணை}}
\textbf{\ta{அரசாணை}} — \emph{aracāṇai} — \textquotedblleft{}a royal decree\textquotedblright{}.

\begin{quote}
\ta{மன்னர்} \ta{ஒரு} \ta{புதிய} \ta{அரசாணையை} \ta{வெளியிட்டார்}.
\end{quote}

\begin{quote}
\emph{maṉṉar oru putiya aracāṇaiyai veḷiyiṭṭār.}
\end{quote}

\begin{quote}
The king issued a new royal decree.
\end{quote}

Today also a government order.

\begin{cousinweb}
\textbf{\ta{அரசாணை}} (\emph{aracāṇai}) joins \ta{அரசு}, \textquotedblleft{}king, government\textquotedblright{}, which DEDR 201 traces to Sanskrit \emph{rājan}, with \ta{ஆணை}, \textquotedblleft{}order, command\textquotedblright{}, whose origin is not settled in the sources at hand.
\end{cousinweb}

\begin{grammarlens}[title={What you've built}]
One more word for telling what happened.
\end{grammarlens}

\subsection*{Your turn}
\begin{itemize}
\raggedright
  \item \emph{Say it:} \emph{aracāṇai}
  \item \emph{Say it:} \emph{aracāṇai}, once more
  \item \emph{Say it:} \emph{maṉṉar oru putiya aracāṇaiyai veḷiyiṭṭār.} — The king issued a new royal decree.
  \item \emph{Recall:} say the Tamil for a royal tribute, then the Tamil for to flourish, then say \emph{aracāṇai} again
\end{itemize}

\begin{tcolorbox}[breakable,colback=teal!4,colframe=teal!35!black,title={Before you move on}]
What does \ta{அரசாணை} mean? (\textquotedblleft{}A royal decree\textquotedblright{}.) Say it once more.
\end{tcolorbox}
\section[\texorpdfstring{\ta{செப்பேடு} (ceppēṭu)}{செப்பேடு (ceppēṭu)}]{\ta{செப்பேடு} (ceppēṭu) — a copper-plate grant}

\label{lesson:TA-C331-ceppetu}



\begin{quote}
Before the new one: say the Tamil for to flourish, then the Tamil for a royal decree.
\end{quote}

\subsection*{You'll want to know: \ta{செப்பேடு}}
\textbf{\ta{செப்பேடு}} — \emph{ceppēṭu} — \textquotedblleft{}a copper-plate grant\textquotedblright{}.

\begin{quote}
\ta{அருங்காட்சியகத்தில்} \ta{பழைய} \ta{செப்பேடுகள்} \ta{இருக்கின்றன}.
\end{quote}

\begin{quote}
\emph{aruṅkāṭciyakattil paḻaiya ceppēṭukaḷ irukkiṉṟaṉa.}
\end{quote}

\begin{quote}
There are old copper-plate grants in the museum.
\end{quote}

Plural: \ta{செப்பேடுகள்}, \emph{ceppēṭukaḷ}.

\ta{செம்பு}, copper + \ta{ஏடு}, leaf: records cut on copper sheets.

\begin{cousinweb}
\textbf{\ta{செப்பேடு}} (\emph{ceppēṭu}) is a Tamil compound built from \ta{செம்பு}, \textquotedblleft{}copper\textquotedblright{} (DEDR 2775), and \ta{ஏடு}, \textquotedblleft{}a strip of palmyra leaf for writing, a book\textquotedblright{} (DEDR 884).
\end{cousinweb}

\begin{grammarlens}[title={What you've built}]
One more word for telling what happened.
\end{grammarlens}

\subsection*{Your turn}
\begin{itemize}
\raggedright
  \item \emph{Say it:} \emph{ceppēṭu}
  \item \emph{Say it:} \emph{ceppēṭu}, once more
  \item \emph{Say it:} \emph{aruṅkāṭciyakattil paḻaiya ceppēṭukaḷ irukkiṉṟaṉa.} — There are old copper-plate grants in the museum.
  \item \emph{Recall:} say the Tamil for to flourish, then the Tamil for a royal decree, then say \emph{ceppēṭu} again
\end{itemize}

\begin{tcolorbox}[breakable,colback=teal!4,colframe=teal!35!black,title={Before you move on}]
What does \ta{செப்பேடு} mean? (\textquotedblleft{}A copper-plate grant\textquotedblright{}.) Say it once more.
\end{tcolorbox}
\section[\texorpdfstring{\ta{வள்ளல்} (vaḷḷal)}{வள்ளல் (vaḷḷal)}]{\ta{வள்ளல்} (vaḷḷal) — a patron; a generous giver}

\label{lesson:TA-C331-vallal}



\begin{quote}
Before the new one: say the Tamil for a royal decree, then the Tamil for a copper-plate grant.
\end{quote}

\subsection*{You'll want to know: \ta{வள்ளல்}}
\textbf{\ta{வள்ளல்}} — \emph{vaḷḷal} — \textquotedblleft{}a patron; a generous giver\textquotedblright{}.

\begin{quote}
\ta{அந்த} \ta{வள்ளல்} \ta{புலவர்களுக்கு} \ta{நிறையப்} \ta{பரிசுகள்} \ta{கொடுத்தார்}.
\end{quote}

\begin{quote}
\emph{anta vaḷḷal pulavarkaḷukku niṟaiyap paricukaḷ koṭuttār.}
\end{quote}

\begin{quote}
That patron gave many gifts to poets.
\end{quote}

A famous figure in old Tamil poems.

\begin{cousinweb}
\textbf{\ta{வள்ளல்}} (\emph{vaḷḷal}) is native Dravidian (DEDR 5304): a person of boundless liberality, from \ta{வள்}, \textquotedblleft{}abundance\textquotedblright{}.
\end{cousinweb}

\begin{grammarlens}[title={What you've built}]
One more word for telling what happened.
\end{grammarlens}

\subsection*{Your turn}
\begin{itemize}
\raggedright
  \item \emph{Say it:} \emph{vaḷḷal}
  \item \emph{Say it:} \emph{vaḷḷal}, once more
  \item \emph{Say it:} \emph{anta vaḷḷal pulavarkaḷukku niṟaiyap paricukaḷ koṭuttār.} — That patron gave many gifts to poets.
  \item \emph{Recall:} say the Tamil for a royal decree, then the Tamil for a copper-plate grant, then say \emph{vaḷḷal} again
\end{itemize}

\begin{tcolorbox}[breakable,colback=teal!4,colframe=teal!35!black,title={Before you move on}]
What does \ta{வள்ளல்} mean? (\textquotedblleft{}A patron; a generous giver\textquotedblright{}.) Say it once more.
\end{tcolorbox}
\section[\texorpdfstring{\ta{கொடை} (koṭai)}{கொடை (koṭai)}]{\ta{கொடை} (koṭai) — a donation; a gift}

\label{lesson:TA-C331-kotai}



\begin{quote}
Before the new one: say the Tamil for a copper-plate grant, then the Tamil for a patron.
\end{quote}

\subsection*{You'll want to know: \ta{கொடை}}
\textbf{\ta{கொடை}} — \emph{koṭai} — \textquotedblleft{}a donation; a gift\textquotedblright{}.

\begin{quote}
\ta{இந்தப்} \ta{பள்ளி} \ta{ஒரு} \ta{வணிகரின்} \ta{கொடை}.
\end{quote}

\begin{quote}
\emph{intap paḷḷi oru vaṇikariṉ koṭai.}
\end{quote}

\begin{quote}
This school is a merchant's donation.
\end{quote}

From \ta{கொடு}, to give.

\begin{cousinweb}
\textbf{\ta{கொடை}} (\emph{koṭai}) is native Dravidian: DEDR 2053 gives it under \ta{கொடு}, \textquotedblleft{}to give\textquotedblright{}, as \textquotedblleft{}giving away as a gift\textquotedblright{}.
\end{cousinweb}

\begin{grammarlens}[title={What you've built}]
One more word for telling what happened.
\end{grammarlens}

\subsection*{Your turn}
\begin{itemize}
\raggedright
  \item \emph{Say it:} \emph{koṭai}
  \item \emph{Say it:} \emph{koṭai}, once more
  \item \emph{Say it:} \emph{intap paḷḷi oru vaṇikariṉ koṭai.} — This school is a merchant's donation.
  \item \emph{Recall:} say the Tamil for a copper-plate grant, then the Tamil for a patron, then say \emph{koṭai} again
\end{itemize}

\begin{tcolorbox}[breakable,colback=teal!4,colframe=teal!35!black,title={Before you move on}]
What does \ta{கொடை} mean? (\textquotedblleft{}A donation; a gift\textquotedblright{}.) Say it once more.
\end{tcolorbox}
\section[\texorpdfstring{\ta{செழி} (ceḻi)}{செழி (ceḻi)}]{\ta{செழி} (ceḻi) — to prosper; to thrive}

\label{lesson:TA-C331-celi}



\begin{quote}
Before the new one: say the Tamil for a patron, then the Tamil for a donation.
\end{quote}

\subsection*{You'll want to know: \ta{செழி}}
\textbf{\ta{செழி}} — \emph{ceḻi} — \textquotedblleft{}to prosper; to thrive\textquotedblright{}.

\begin{quote}
\ta{ஆற்றின்} \ta{கரையில்} \ta{நகரங்கள்} \ta{செழித்தன}.
\end{quote}

\begin{quote}
\emph{āṟṟiṉ karaiyil nakaraṅkaḷ ceḻittaṉa.}
\end{quote}

\begin{quote}
Towns prospered on the banks of the river.
\end{quote}

Past (it …): \ta{செழித்தது}, \emph{ceḻittatu}.

Of plants too: to grow well.

\begin{cousinweb}
\textbf{\ta{செழி}} (\emph{ceḻi}) is native Dravidian (DEDR 2789): to thrive, flourish, prosper, as a crop or a kingdom.
\end{cousinweb}

\begin{grammarlens}[title={What you've built}]
One more word for telling what happened.
\end{grammarlens}

\subsection*{Your turn}
\begin{itemize}
\raggedright
  \item \emph{Say it:} \emph{ceḻi}
  \item \emph{Say it:} \emph{ceḻi}, once more
  \item \emph{Say it:} \emph{āṟṟiṉ karaiyil nakaraṅkaḷ ceḻittaṉa.} — Towns prospered on the banks of the river.
  \item \emph{Recall:} say the Tamil for a patron, then the Tamil for a donation, then say \emph{ceḻi} again
\end{itemize}

\begin{tcolorbox}[breakable,colback=teal!4,colframe=teal!35!black,title={Before you move on}]
What does \ta{செழி} mean? (\textquotedblleft{}To prosper; to thrive\textquotedblright{}.) Say it once more.
\end{tcolorbox}
